% Supplementary Information of the ERL Letter (D-107): copy of methods_datums.tex, adapted (ICESat-2
% through EIGEN-6C4, D-078; GLO-30 and GEDI, D-091; the BHS-77 offset, D-093).
% Methods: the vertical datum chain (docs/DATUMS.md). Design: D-069, D-070, D-073, D-074;
% PREREG section 5 and Amendment 11. Every number comes from \metric{...}.

\subsection*{Vertical datums}
\emph{Working datum.} All heights were compared on the reference of the DAHITI series: normal heights, ellipsoidal altimeter heights minus a quasi-geoid model\citep{schwatke2015}, EIGEN-6C4 in the current data version\citep{dahiti_wl,forste2014}. Level bins, waterlines, topography and hindcast all use it. Levels are reported in the Baltic Height System; a constant offset between the two systems applies equally to levels and to heights brought to the DAHITI reference, so it cancels in every comparison here and is estimated only for reporting (pre-registration section~5).

\emph{Relief grids.} For the total error of the hindcast audit, the grids were read as heights above the EGM2008 geoid\citep{pavlis2012}. ETOPO~2022 states this reference in its file metadata\citep{macferrin2025}; GEBCO and ETOPO1 describe their heights as referred to sea level\citep{gebco2026,amante2009}, and the SRTM15+ distribution files state none\citep{tozer2019}, so for these grids the EGM2008 reading is an assumption. Each grid was brought to the DAHITI reference by adding, cell by cell, $N_{\mathrm{EGM2008}} - N_{\mathrm{EIGEN\text{-}6C4}}$, computed from the two tide-free spherical-harmonic models with the calculation service of the International Centre for Global Earth Models\citep{ince2019}. Over the domain the difference ranges from \metric{datum.dN_min_m} to \metric{datum.dN_max_m}~m (standard deviation \metric{datum.dN_sd_m}~m). The widely distributed realisation of EGM2008, which ETOPO~2022's geoid layer reproduces, carries a further constant of \metric{datum.nga_offset_m}~m\todo{cite: NGA EGM2008 distribution documentation}; the total error uses it, and the total error without it is a sensitivity. The shape error fits its own constant and depends on neither choice.

\emph{Lidar.} ICESat-2 ATL06 and ATL08 heights (Release~007) are tide-free ellipsoidal heights on WGS84 in the ITRF2020 frame\citep{atl06_v7,atl08_v7,is2_compguide_r007}. Heights above the geoid of the working datum are $H = h - N_{\mathrm{EIGEN\text{-}6C4}}$, with $N$ interpolated bilinearly at each point from the tide-free grid of the same calculation service, for ATL06 and ATL08 alike; the EGM2008 geoid height that ATL06 carries is used only as a datum sensitivity. A constant $c$, the median of the topography minus $H$ over training-block points, with a block-bootstrap uncertainty, then absorbs the remaining constants (zero-degree term, frame and tide-system offsets): \metric{h2.atl06.datum_constant_m}~m ($\sigma$ \metric{h2.atl06.datum_constant_sigma_m}~m) for ATL06 and \metric{h2.atl08.datum_constant_m}~m ($\sigma$ \metric{h2.atl08.datum_constant_sigma_m}~m) for ATL08. Both $h$ and $N$ are tide-free, so no tide conversion is applied. DGFI does not state the tide system of DAHITI heights; if it is mean-tide, a north--south tilt of at most a few centimetres remains and is stated as an error. GEDI heights (pre-registration Amendment~2) are converted as ICESat-2 and get their own training-block constant. GLO-30 heights are referenced to EGM2008\citep{copdem_handbook_i5} and were brought to the working datum with the per-cell difference used for the relief grids and their own training-block constant; its two geocells in the domain without TanDEM-X measurements were excluded.

\emph{Reporting datum.} Levels are reported in the Baltic Height System of 1977 (BHS-77) through the offset $\delta = H_{\mathrm{BHS}} - H_{\mathrm{DAHITI}}$, estimated from the official monthly levels of the CASPCOM General Catalogue\citep{caspcom2026} at the \metric{datum.bhs.n_posts} ``century'' posts from which CASPCOM forms its sea-mean level, over the \metric{datum.bhs.n_years} complete years between \metric{datum.bhs.first_year} and \metric{datum.bhs.last_year} that both series cover (\metric{datum.bhs.n_post_years} post-years). Per year and post, the annual mean of the post minus the annual mean of DAHITI was averaged over years for each post and then over posts: $\delta = \metric{datum.bhs.delta_m}$~m, with a standard error of \metric{datum.bhs.sigma_years_m}~m from a moving-block bootstrap over years and \metric{datum.bhs.sigma_posts_years_m}~m when posts are resampled as well; the latter is the uncertainty quoted, because the per-post offsets range from \metric{datum.bhs.post_min_m} to \metric{datum.bhs.post_max_m}~m. Sensitivities give \metric{datum.bhs.sens_min_m} to \metric{datum.bhs.sens_max_m}~m. Because $\delta$ is added equally to levels and heights, it cancels in every comparison of this study.

\emph{SWOT.} SWOT elevations are heights above a geoid that the files do not name, after removal of solid-earth, load and pole tides. Its difference from the DAHITI reference is constant in time within a block and is removed with each block's mean departure in the setup model.
